\section{The double-cover displacement and its elliptic tails}\label{sec:smooth-crossings}
We prove the double-cover displacement used in \S\ref{p10-sec:incidence} and separate its first contracted elliptic contribution into the moving and ramification-point terms.
\subsection{The local contribution at a double-cover wall}\label{p9-sec:wall}

Let $C$ be a smooth curve of genus $g$, and let $N\to C$ be a complex
bundle of rank two and degree $2g-2$. Consider a smooth path $D_t$ of
real Cauchy--Riemann operators on $N$, with $D_0$ invertible. We orient
the parameter interval by increasing $t$. Write $\mathcal H$ for the
space of smooth connected double covers of the fixed target, locally
taken in a uniformizing chart. At a cover $f:Y\to C$ with $b$ branch
values, these values are local complex coordinates up to ordering.
The deck involution acts trivially on these coordinates.

Suppose $K=\ker f^*D_0$ is a real line, and put $P=\operatorname{coker}
f^*D_0$. Differentiating the family of pullback operators on a section
in $K$ and projecting to $P$ gives the intrinsic derivative
\begin{equation}\label{p9-eq:wall-differential}
 W_f:T_f\mathcal H\oplus\RR_t\longrightarrow
                  \operatorname{Hom}_{\RR}(K,P).
\end{equation}
We use the complex orientation of $\mathcal H$ and the absolute
Cauchy--Riemann determinant orientation of $f^*D_0$. Indeed,
Riemann--Hurwitz gives $\dim_{\RR}P=2b+1$, and the determinant of the
space on the right is $(K^*)^{\otimes 2b}\otimes(\det f^*D_0)^*$.
Write $\eta_f$ for the sign of \eqref{p9-eq:wall-differential} when it is an
isomorphism. This sign uses the full pullback operator, including its
invariant summand \cite[(4.59) and (5.59)--(5.60)]{bai}.

\begin{theorem}\label{p9-thm:wall}
Fix an arithmetic genus bound $k_*$. Suppose the smooth connected
double covers of genus at most $k_*$ with nonzero pullback kernel at
$t=0$ form a finite set up to isomorphism. For each such $f$, assume
that the kernel is an anti-invariant real line, the map
$(f,\xi):Y\to N$ is an embedding for a generator $\xi$, and
\eqref{p9-eq:wall-differential} is an isomorphism. For sufficiently small
$\delta>0$ the pullback operators at the two endpoints are injective on every
degree-two stable map of genus at most $k_*$. In each
genus $k\le k_*$ the change of the absolute normal Euler number is
\begin{equation}\label{p9-eq:wall-jump}
 L_{2,k}(D_{\delta})-L_{2,k}(D_{-\delta})
 =\sum_{f:Y\to C}\eta_f\,[v^{2k-2}]
                \bigl(2\sin(v/2)\bigr)^{2g(Y)-2}.
\end{equation}
Each summand is the change supported on stable maps that factor through a degree-one
stable map to the source of the indicated cover. In particular a single wall with smooth
source of genus $H$ contributes the series
$\eta_f(2\sin(v/2))^{2H-2}$ coefficientwise.
\end{theorem}

\begin{proof}
We first treat a single cover $f$ of genus $H$. Choose a complex
connection on $N$. On a disk bundle about the zero section we realize
$D_t$ as the normal operator of an almost complex structure with
holomorphic projection and standard complex structure on each fiber.
For a holomorphic map $h$ to $C$, prescribe the equation of a lift
$(h,u)$ by
\begin{equation}\label{p9-eq:wall-cubic}
 h^*D_t u+h^*B(u,u,u)=0 ,
\end{equation}
where $B$ is a smooth real symmetric trilinear form on $N$ with
values in $\Omega_C^{0,1}\otimes N$. Here pullback of the form includes
composition with $dh$. In a connection splitting, the difference of
$D_t$ from the connection Dolbeault operator and the cubic both specify
vertical-valued $(0,1)$-forms in the horizontal variable. Such a form
specifies the off-diagonal block of the almost complex structure, giving
\eqref{p9-eq:wall-cubic} for every $h$ \cite[Lemma~5.9]{bai}. We multiply
the change of horizontal structure by a radial cutoff equal to one in
the disk under consideration and extend across $\PP(N\oplus\CC)$.
A unitary connection gives a closed form restricting to the
Fubini--Study form on the fibers. Adding a sufficiently large multiple
of a positive area form from $C$ tames the extended family, since the
horizontal terms are bounded on the compact bundle.

\smallskip
\noindent\emph{The smooth branch and its sign.}
The invariant part of $f^*D_0$ has index zero. To see this, deform
equivariantly to a pullback Dolbeault operator and use
$(f_*\mathcal O_Y)^+=\mathcal O_C$, together with exactness of
invariants under the deck group in characteristic zero. The hypothesis
on $K$ thus implies that $P$ is entirely anti-invariant. Take
equivariant complements to the kernel and image in a Sobolev completion
of \eqref{p9-eq:wall-cubic} over the Hurwitz chart. The implicit function
theorem solves the image equation for the kernel complement. In local
coordinates $y$ on the chart and with kernel coordinate $a\xi$, the
remaining equation is
\[
                    a F(y,t,a)=0,\qquad F(y,t,-a)=F(y,t,a).
\]
Indeed the zero section solves the equation at all $(y,t)$ and the
deck change negates both $a$ and the obstruction in $P$. The
differential of $F$ in $(y,t)$ at the origin equals $W_f$ evaluated
on $\xi$. Its coefficient of $a^2$ at $(y,t)=0$ is the image of
$f^*B(\xi,\xi,\xi)$ in $P$. The solution in the kernel complement
starts in degree three there, and its linear image projects to zero.

Let $p$ be a nonzero adjoint section annihilating the image of the
Hurwitz directions in $P$. On an unramified disk where $p$ and $\xi$
are nonzero on one sheet, prescribe an arbitrary compactly supported
form on that sheet as $f^*B(\xi,\xi,\xi)$. This uses a real covector
taking value one on $\xi$ and its third tensor power. Such a disk
exists by unique continuation. On the other sheet both the prescribed
form and the adjoint change sign, so their two integrals agree.
Choose the form to give a nonzero pairing. The implicit function
theorem for $F=0$ now gives even functions $y(a),t(a)$ with
\[
                     t(a)=c a^2+O(a^4),\qquad c\ne0 .
\]
The two signs of $a$ give the same curve by the deck action. Dividing
its vertical coordinate by $a$ gives smooth convergence to
$(f,\xi)$, so the curve for small $a>0$ is embedded and has trivial
stabilizer.

At a nonzero solution the derivative of $aF$ in $(y,a)$ has first
columns $aF_y$. Its last column modulo these is $-a t'(a)F_t$ by
differentiation along the branch. It is invertible, and comparison
with $W_f$ gives the sign $-\operatorname{sgn}(t)\eta_f$. In this
comparison the tangential deformation of the cover has the complex
Hurwitz orientation. Indeed its holomorphic deformation complex reduces
to the torsion cokernel of $df$, supported at the simple ramification
points. At the zero section the vertical and tangential determinants
split, and continuation gives the determinant of the embedded graph.
This is the absolute sign convention of
\cite[Theorem~4.12 and (4.30), (4.59)]{bai}. Thus the change in the
signed number of embedded curves with nonzero vertical coordinate is
$-\eta_f$.

\smallskip
\noindent\emph{Isolation through genus $k_*$.}
Suppose there were other simple curves, of degree at most two over $C$
and normalization genus at most $k_*$, in arbitrarily small closed
tubes at parameters tending to zero. Pass to a sequence with fixed
normalization genus. There is no nonconstant compact curve in an
affine fiber. In degree one the projection of the normalization is an
isomorphism, and invertibility of $D_0$ excludes a small nonzero graph.
Equivalently, a sequence of such graphs can be divided by its maximum
norm, and elliptic compactness over the fixed base would give a
nonzero element of $\ker D_0$.

In degree two write the maps from the smooth normalizations as
$(h_n,u_n)$ and put $\rho_n=\max |u_n|>0$. The graphs
$(h_n,u_n/\rho_n)$ solve \eqref{p9-eq:wall-cubic} with cubic coefficient
$\rho_n^2B$. On a disk bundle of radius greater than one these
structures extend to a smooth two-parameter family in $(t,\rho)$
including $(0,0)$. A common taming form as above has bounded area on
the graphs, since its cohomology class on the disk bundle is pulled
back from $C$ and the projection degree is fixed. Gromov compactness,
with a point marked where the norm is one, gives a stable limit in
the linear structure of $D_0$ whose vertical coordinate is nonzero.
Forget the extra mark and stabilize. If two normalized components
had degree one, invertibility of $D_0$ would force both sections to
vanish and all constant components would have the same zero vertical
coordinate at their attachments. Thus the unique effective
normalization is a smooth double cover with a nonzero kernel section.
By hypothesis its graph is an embedded curve $\Gamma=(f,\lambda\xi)(Y)$,
with $\lambda\ne0$.

The stable limit factors in degree one through $\Gamma$. Its effective
normalization maps isomorphically to $\Gamma$, so no self-node
identifies two points on that normalization. Each maximal connected
contracted subcurve attaches at a single point, since its image has a
single preimage on the effective normalization. For the resulting
map $w$ to $\Gamma$ and every real normal Cauchy--Riemann operator
$D_\Gamma$, restriction and pullback give inverse isomorphisms
\begin{equation}\label{p9-eq:wall-degree-one-ker}
                     \ker w^*D_\Gamma \simeq \ker D_\Gamma .
\end{equation}
A section on a contracted component is holomorphic with values in a
fixed vector space, hence constant, and the node conditions prescribe
its value on the whole contracted subcurve. The compactness theorem
for sequences of simple maps to a finite-dimensional family of almost
complex structures \cite[Theorem~3.1 and Remark~3.2(b)]{bai} says that
a limit over an embedded curve, unless the limiting map is an
isomorphism to that curve, has a normal kernel section outside the
pullbacks from the embedded curve. By \eqref{p9-eq:wall-degree-one-ker}
the limit is an isomorphism. We thus have smooth convergence modulo
parametrization. The original graphs eventually belong to the
Hurwitz chart and its implicit-function neighborhood, contradicting
their choice.

We also need injectivity at the endpoints on the entire fixed-target
space in genus $k\le k_*$. Apply the same compactness argument to
graphs of normalized kernel sections of the linear operators at
$t_n\to0$. The invariant part of any smooth double-cover kernel
descends to a base section. It descends on the unramified locus by
invariance and extends across the branch values by bounded
removability for the base equation. Since $D_t$ remains invertible,
the kernel sections are odd. Their nonzero graphs are simple because
the two sheets have different values wherever the section is nonzero.
After smooth convergence in the rescaling argument, the linear problem
in the Hurwitz chart reduces to $F(y,t,0)=0$, whose only solution
near the origin is $(0,0)$ by \eqref{p9-eq:wall-differential}.
Thus all smooth covers through the cutoff have injective pullback for
small $0<|t|$. For a nodal source, restriction to its effective
normalizations and the constant matching on its contracted components
give injectivity as well. At $t=0$ the same reasoning shows that every
stable map with kernel lifts in degree one to $Y$ via its graph in
$\Gamma$. Conversely every such lift carries the pullback of $\xi$.

\smallskip
\noindent\emph{Conservation of the local degree.}
The isolation just proved gives a closed tube of some radius $r>0$
and an interval about zero in which the only simple curves in that
tube through the cutoff are $C$ and the embedded curves found above.
Shrink the interval so the latter lie in the tube of radius $r/2$.
Every effective component of a degree-two stable map contained in the
closed tube factors through a simple normalization of degree at most
two and genus at most $k_*$. Connectedness therefore prevents any
such map from meeting the boundary. In the compactified target these
maps form a compact open and closed part of the parametrized stable-map
space. Compactness gives closedness, and support in the interior gives
openness. Its virtual degrees at the two ends agree by fixed-genus deformation invariance and the splitting of the virtual
class over open and closed substacks \cite[\S3.1, (3.4)--(3.5)]{diw}.

At an endpoint the maps supported on $C$ contribute $L_{2,k}(D_t)$
by the normal Euler comparison \cite[Theorem~1.2, (1.7)]{zinger}.
The maps on the additional embedded curve, when present, have degree
one over it. We compute their degree including its orientation.
Let $D'$ be the invertible normal operator on that embedded curve
of genus $H$. For any deformation of $D'$ on the curve, the pullback and base operators
admit compatible finite-dimensional stabilizations. The quotient of the
two resulting complexes has zero kernel on every degree-one stable
map $w$. Indeed
\eqref{p9-eq:wall-degree-one-ker} identifies the kernels, and restriction
to the effective normalization is a left inverse on cokernels to the
pullback of base forms. The quotient cokernel is therefore a vector
bundle throughout the deformation, by the stabilized construction in
\cite[Appendix~B, (B.8)--(B.10)]{bai}. Its determinant is the
quotient of the pullback and base determinants. At an invertible
endpoint the Euler class in the absolute orientation is the Euler
class of this bundle multiplied by $\epsilon(D')$.

The normal bundle has degree $2H-2$, since the ambient first Chern
class on the tube vanishes. Deform $D'$ to the Dolbeault operator of
two line bundles of degree $H-1$ off the theta divisor, using
$\mathcal O(-1)^{\oplus2}$ if $H=0$. There the cokernel has its
complex orientation and its virtual Euler degree is the ordinary
degree-one local invariant. Pandharipande's degree-one formula
\cite[Theorem~2]{pandhodge} gives, in our genus variable, the series
$(2\sin(v/2))^{2H-2}$. This also gives the degree of the reduced
cokernel at $D'$ by deformation. The change from the embedded curves
off $C$ is consequently $-\eta_f$ times this series through the
cutoff. Conservation gives \eqref{p9-eq:wall-jump} with the opposite
sign on $C$.

\smallskip
\noindent\emph{Simultaneous covers.}
For finitely many covers $f_i$ a single cubic $B$ can be chosen
outside the finite union of hyperplanes in the adjoint tests above.
The compactness argument then gives one branch in each indicated
Hurwitz chart and the sum in \eqref{p9-eq:wall-jump}. To identify the
individual supported changes, we separate their times. With generators
$\xi_i$ and adjoints $p_i$ as above, the conormal to the $i$th wall
on real zeroth-order variations of the base operator is
\[
                   \ell_i(A)=\langle p_i,f_i^*A\,\xi_i\rangle .
\]
For distinct covers these functionals are nonproportional. Work on a
common unramified disk where the sections and adjoints in question
are nonzero. The two sheet integrals for each $\ell_i$ agree. Testing
with arbitrary real-linear $A$ on this disk, proportionality would
therefore imply proportionality of the rank-one tensor densities,
and in particular real proportionality $\xi_i=c\xi_j$ on chosen
sheets. The two sections solve the same equation there, so
$\bar\partial c=0$ and $c$ is a nonzero constant. Off the union of
the branch values the two covers have distinct nonzero mod-two
characters, since a cover of the complement has a unique finite
normal completion. Their characters are independent and their fiber
product is connected. Unique continuation extends the equality over
this fiber product. A deck transformation flipping only the first
sheet then negates exactly one side, a contradiction.

Choose $\chi(t)$ equal to one near zero and supported away from the
endpoints. Write $t_i(s)$ for the crossing time of $f_i$ in the
family $D_t+s\chi(t)A$. Then
\[
              t_i'(0)=-\ell_i(A)/\ell_i(\dot D_0).
\]
The denominators
are nonzero by \eqref{p9-eq:wall-differential}. Pairwise
nonproportionality lets us choose $A$ so all derivatives are distinct.
The compactness argument for the linear graphs, now including $s$,
confines all kernels through the cutoff to the same charts. Thus the
perturbation separates the crossings and preserves their signs.

For completeness, let $M=\overline{\mathcal M}_{k}(C,2)$, let $I$
be the time interval, and let $Z$ be the closed noninjectivity locus
in $M\times I$. On its complement the absolutely oriented pullback
cokernel has an Euler class $e$. The image of $e$ under the
connecting morphism of the pair $(M\times I,(M\times I)\setminus Z)$
pairs with $[M]^{\rm vir}\times[I,\partial I]$ to give the final
Euler degree minus the initial one. At $s=0$ the kernel supports of
the distinct covers are disjoint compact images of
$\overline{\mathcal M}_k(Y_i,1)$, since the effective normalization
recovers $f_i$. Choose disjoint neighborhoods. For sufficiently small
$s$ all kernels remain in them and avoid their boundaries, by
properness and the local stabilized operators
\cite[Appendix~B, (B.8)]{bai}. The relative class for the whole
$(t,s)$-family splits over these neighborhoods by excision. Pairing
over the product virtual family shows that each summand has the same
degree at $s=0$ as at small positive $s$. There each is the
single-cover contribution already computed. All compactness choices
were made at the prescribed finite cutoff.
\end{proof}

On a fixed-target path with finitely many degree-at-most-two wall
times in all genera, the coefficients in Theorem~\ref{p9-thm:wall} may
be evaluated at common regular endpoints on either side of a single
time. On a regular parameter interval the pullback cokernels form an
oriented bundle on each compact fixed-genus map space, so their
Euler degrees are constant there. For generic fixed-target paths the
finiteness assertion follows from \cite[Definition~2.6,
Proposition~2.7 and Theorem~3.40]{dw}.
\subsection{A contracted elliptic tail at a ramification point}\label{p9-sec:tail}

Let $C$ be a smooth curve of genus $g\geq 2$, let $N\to C$ be a complex
bundle of rank two and degree $2g-2$, and let $D_t$ be a smooth path of
real Cauchy--Riemann operators on $N$. Fix a smooth connected cover
$f:Y\to C$ of degree $m$ and genus $H$, with $b$ simple ramification
points over distinct values. Write $G=\operatorname{Aut}_C(f)$ and
$a=|G|$. The fixed-target Hurwitz stack $\mathcal H$ is smooth of complex
dimension $b$ at $f$. Indeed, the quotient of $T_Y\to f^*T_C$ is a
skyscraper sheaf of length $b$, with one summand at each ramification point.

Suppose $K=\ker f^*D_0$ is a real line and the intrinsic derivative
\begin{equation}\label{p9-eq:tail-core}
 T_f\mathcal H\oplus\RR_t
 \longrightarrow \operatorname{Hom}_{\RR}(K,\operatorname{coker} f^*D_0)
\end{equation}
is an isomorphism. Its sign $\eta$ uses the complex orientation of
$\mathcal H$, increasing time and the absolute determinant orientation
transported from complex-linear operators. For $0\ne\xi\in K$ assume
$d\xi_r\ne 0$ at each ramification point $r$. In a frame pulled back from
the target, $d\xi_r:T_rY\to N_{f(r)}$ is independent of the connection
because $df_r=0$, and is complex-linear by the pullback equation.

\begin{theorem}\label{p9-thm:tail}
In source genus $k=H+1$, the stable maps obtained by attaching a
contracted one-pointed curve of arithmetic genus one to $f$ form an
isolated compact part of the locus with nonzero normal kernel at the
crossing. The local degree of the change of the absolute normal Euler class
supported on this family is
\begin{equation}\label{p9-eq:tail-total}
 -\frac{2\eta}{a}\left(\frac{m(g-1)}{12}+\frac{b}{24}\right)
 =-\frac{2\eta}{a}\frac{H-1}{12}.
\end{equation}
The first term comes from the component on which the tail persists.
The second is the sum over ramification points on the component on
which the attaching node smooths. The corresponding local degree at the smooth cover in genus $H$ is
$2\eta/a$. When the endpoint pullbacks are injective in the genus under
consideration, each local degree is a summand of the difference of the
endpoint Euler numbers.
\end{theorem}

\begin{proof}
For fixed $k$ let $X=\overline{\mathcal M}_k(C,m)$ be the stack of
connected stable maps, let $I$ be the parameter interval, and let
$U\subset X\times I$ be the open set where the pullback normal operator
is injective. Its oriented real cokernel bundle $Q_U$ has rank
$d=4(k-1-m(g-1))$, the real virtual dimension of $X$. The relative class
\begin{equation}\label{p9-eq:tail-half}
 u=\tfrac12\delta e(Q_U)\in H^{d+1}(X\times I,U;\QQ)
\end{equation}
defines the local degrees in the statement. Here $\delta$ is the
connecting homomorphism. With injective endpoints,
$2u$ paired with the image of $[X]^{\rm vir}\times[I,\partial I]$
is the difference of the endpoint Euler numbers by the boundary formula.
Excision gives the summand of this pairing supported near an isolated
compact part of the complement of $U$. This local pairing also makes
sense when other covers have kernel at the same time.

Near a map with kernel line $K$ and cokernel $Q$, stabilization gives
a finite-dimensional presentation of the normal operator
\cite[\S6.1 and Appendix~B, (B.8)]{bai}. In its matrix space the
stratum with kernel line has normal space $\operatorname{Hom}_{\RR}(K,Q)$.
Its determinant is $(K^*)^{\otimes d}\otimes(\det D)^*$, where $\det D$
denotes the Fredholm determinant at that map. Thus the absolute
orientation coorients this stratum. On the link of a normal disc,
choosing a generator of $K$ identifies the universal cokernel with
the tangent bundle of $S^d$. Its Euler degree is two with the boundary
orientation. For $d=0$ the two invertible matrices have opposite
determinant signs. Hence \eqref{p9-eq:tail-half} pulls back the positive
Thom class of this stratum on the corank-at-most-one neighborhood.
Two stabilizations admit a common sum, which adds an identity block
after a change of splitting. The normal space and its orientation are
preserved. In particular \eqref{p9-eq:tail-core} gives local degree
$\eta/a$ for $u$ at the smooth cover.

Write $(E,e)$ for the pointed tail, $p$ for the attaching point on $Y$,
and $\mathbb E$ for the Hodge line on $\overline{\mathcal M}_{1,1}$.
The compact family over $(f,0)$ is
\[
 T=[(Y\times\overline{\mathcal M}_{1,1})/G].
\]
We describe the fixed-target virtual class near $T$. The cover
deforms unobstructedly, and evaluation of the constant value on $E$
eliminates the matching condition at the node. The remaining obstruction
space is $T_{f(p)}C\otimes H^1(E,\mathcal O_E)$. Parameters consist of
the cover, the point $p$, the pointed tail, and the attaching parameter
$q\in T_pY\otimes T_eE$. The node-preserving maps exist for every value
of the first three parameters, so their scalar holomorphic obstruction
equation has the form $q F=0$. On $q=0$ its derivative in the smoothing
direction is $df_p\otimes\delta_e$, where
$\delta_e:T_eE\to H^1(E,\mathcal O_E)$ is the boundary map for
$\mathcal O_E(e)$. Indeed, substitution of the plumbing relation into
the linear term of the core map gives the simple principal part on $E$.
The dual of $\delta_e$ is evaluation of a dualizing differential at the
mark, an isomorphism also on the irreducible nodal elliptic curve.
Thus $F$ is a unit away from ramification. At a simple ramification
point it vanishes simply as $p$ varies. Locally there are two smooth branches
$B=\{q=0\}$ and $M=\{F=0\}$ near such a point. The fixed-target
virtual class is the divisor of the scalar obstruction section and hence
the sum of their fundamental cycles, each with multiplicity one. At a
nodal tail the modulus of the pointed elliptic curve includes the
independent smoothing parameter of its self-node.

On $B$, normalizing at the attaching node identifies the normal kernel
with the core kernel. Constant tail sections evaluate onto the matching
fiber, so the cokernel over $T$ is the sum of the core cokernel and the
complex two-plane $V=f^*N\otimes\mathbb E^\vee$. Taking the Thom class
along \eqref{p9-eq:tail-core} leaves the Euler class of
$\operatorname{Hom}_{\RR}(K,V)$. On the cover
$Y\times\overline{\mathcal M}_{1,1}$, $K$ is constant and either
generator induces the complex orientation on this excess bundle. The
degree on $B$ is therefore
\begin{equation}\label{p9-eq:tail-moving}
 \frac{\eta}{a}\int_{Y\times\overline{\mathcal M}_{1,1}}
 c_2(f^*N\otimes\mathbb E^\vee)
 =-\frac{\eta}{a}\frac{2m(g-1)}{24}.
\end{equation}
Indeed $c_2=-c_1(f^*N)c_1(\mathbb E)$ and
$\int_{\overline{\mathcal M}_{1,1}}c_1(\mathbb E)=1/24$.
The latter degree follows from the Weierstrass presentation as the
weighted projective stack of weights $4,6$, with Hodge line
$\mathcal O(1)$. Transfer from the
cover accounts for the action of $G$, including its action on $K$.

It remains to restrict $u$ to $M$ at a ramification point $r$. Work on
a uniformizing chart before division by $G$. Choose core and target
coordinates so that $w=w_0+z^2$. On the pointed elliptic family choose
a meromorphic function $P$ with its sole pole of order two at $e$.
This choice extends locally across the nodal fiber by base change for
$\mathcal O_E(2e)$, which has two sections, with only the constants
in $H^0(E,\mathcal O_E(e))$. Take a coordinate $y$ with $P=y^{-2}$
near $e$. The plumbing
\begin{equation}\label{p9-eq:tail-plumbing}
 zy=q,\qquad
 w=w_0+z^2\quad\text{on the core},\qquad
 w=w_0+q^2P\quad\text{on the tail}
\end{equation}
gives maps agreeing on the overlap. Varying the marked ramification
cover and the elliptic modulus gives every tangent to $M\cap B$, and $q$
smooths the attaching node to first order. The inverse function theorem
on the smooth branch $M$ proves that these are local parameters there.

Use cokernel representatives supported on the core outside the attaching
disc and on the tail away from both nodes. The former span by unique
continuation for the adjoint equation, the latter by duality with the
Hodge line. Stabilization by these spaces is surjective with kernel a
real line near the central map. Normalize its section $s$ by a real
evaluation on the thick part of the core, giving $\xi$ at the center.
The stabilized kernels are continuous across the nodal fibers
\cite[Appendix~B]{bai}. In the plumbing norms, a continuous local
section is uniformly bounded and converges on compact subsets away from
the nodes \cite[\S\S4.3--4.4, Proposition~4.15]{siebert}. For the
pulled-back zeroth-order term the multiplication norm is bounded by the
sup norm of the section times the weighted $L^p$ norm of the map
derivative. A map holomorphic on the total space has bounded derivative
on each half-collar in its larger coordinate, so this latter norm tends
to zero on shrinking collars. On the remaining compact sets the
coefficients vary smoothly. The elliptic
obstruction coefficient $c_E$ vanishes identically along $q=0$, where
the stabilized core section extends by its constant value to the tail.

We compute $c_E$ to first order. In a holomorphic frame for a reference
Dolbeault operator on the target disc, the equation on the truncated
elliptic side $|y|\geq |q|/\rho$ reads
\[
 \bar\partial s+A_t(w)(s)\,\overline{dw}=\theta_E(c_E).
\]
Here $A_t(w)$ is real-linear on the fibers. Normalize the forcing
$\theta_E$ by pairing with a regular differential
$\omega=(a_0+O(y))\,dy$ near $e$. With the trace convention for
$\delta_e$, the boundary integral in Stokes' formula becomes
\[
 q a_0\,\frac{1}{2\pi i}\oint_{|z|=\rho}s(z)\frac{dz}{z^2}
       + O_\rho(|q|^2).
\]
For fixed $\rho$ the contour coefficient tends to that of $\xi$ as all
parameters approach the chosen point of $T$. Since $\bar\partial\xi(r)=0$,
it is $\partial_z\xi(0)+O(\rho)$ in the limit. For fixed small $c>0$,
the zeroth-order term on $|q|/\rho\leq |y|\leq c$ is bounded by
\[
 C|q|^2\int_{|q|/\rho}^c u^{-2}\,du\leq C\rho |q|.
\]
Outside this disc $dw=q^2dP$ and the integral of $|dP|\,|\omega|$ is
uniformly bounded, also on a smoothing $uv=\tau$ of the self-node.
Indeed $\omega$ is a bounded multiple of $du/u$, and the derivative
of the regular function $P(u,\tau/u)$ is bounded on $|u|\geq
\sqrt{|\tau|}$. The other half is covered by $v$. On the nodal fiber the two
residues in Stokes' formula cancel by matching of $s$ on the normalization.

Dividing by $q$, then sending $\rho$ to zero, gives the elliptic
coefficient to order $o(|q|)$ as all parameters tend to the wall. Its
linear term in the attaching direction is the intrinsic homomorphism
\begin{equation}\label{p9-eq:tail-jet}
 \begin{aligned}
 T_rY\otimes T_eE &\longrightarrow
 \operatorname{Hom}_{\RR}\bigl(K,N_{f(r)}\otimes H^1(E,\mathcal O_E)\bigr),\\
 v\otimes v'&\longmapsto
 \bigl[\xi\longmapsto d\xi_r(v)\otimes\delta_e(v')\bigr].
 \end{aligned}
\end{equation}
This limit holds along arbitrary simultaneous variations of the cover,
time and elliptic modulus. Compactness of $\overline{\mathcal M}_{1,1}$
gives a uniform estimate in finitely many charts. The map
\eqref{p9-eq:tail-jet} embeds the complex smoothing line since $d\xi_r\ne0$.
In a tubular neighborhood extend the splitting of the obstruction space
from the divisor. The straight-line homotopy from the elliptic
coefficient to the extended linear term stays nonzero for $q\ne0$
after shrinking the neighborhood. With this coefficient fixed, retract
the core coefficient to its restriction at $q=0$, where its derivative
is \eqref{p9-eq:tail-core}. These homotopies require only continuity of the
core coefficient in $q$. They apply over the whole ramification-marked
family: the leading map transforms as the tensor in \eqref{p9-eq:tail-jet},
and splittings can be chosen equivariantly using a metric and common
stabilizations.

After these two Thom isomorphisms the excess bundle on the deck cover
is the complex line
\[
 \mathbb E^\vee\otimes
 \bigl(N_{f(r)}/d\xi_r(T_rY)\bigr).
\]
Its degree is $-1/24$. The smoothing line has its complex orientation.
Interchanging it with the oriented real normal space
to the wall of smooth covers preserves the sign.
Thus the degree of $u$ on $M$ for the orbit of $r$ is
$-\eta/(24|\operatorname{Aut}_C(f,r)|)$. Summation over the orbits
gives $-\eta b/(24a)$. The transversality on $B$ and the nonvanishing
for $q\ne0$ on $M$ isolate $T$ in the kernel locus. Adding the two
virtual-cycle contributions, doubling by \eqref{p9-eq:tail-half}, and
using $2H-2=2m(g-1)+b$ proves \eqref{p9-eq:tail-total}.
\end{proof}

For a double cover of a genus-two target with two branch points,
$H=4$ and $a=2$. The smooth-cover jump is $\eta$. In genus five
the moving attachment contributes $-\eta/6$, and smoothing at either
ramification point contributes $-\eta/24$, giving $-\eta/4$ in all.
In general the ratio $-(H-1)/12$ of the two successive jumps is the
quadratic coefficient of $(\sin(v/2)/(v/2))^{2H-2}$, as in the
degree-one formula \cite[Theorem~2]{pandhodge}.
